% The next line makes it so it compiles the main file if I hit the compile button
% when editing this file in TeXworks.
% !TEX root = ../DiscreteMathAndAlgorithmsBook.tex

\newpage

\section{Problems}

\begin{pro}
Give the degrees of the vertices of each of the following graphs.
Assume $m$ and $n$ are positive integers.
For instance, for $P_n$, $n-1$ of the vertices have degree $2$, and
$2$ vertices have degree $1$.
\begin{lenum}
\item
$C_n$
\item
$Q_n$
\item
$K_n$
\item
$K_{m,n}$
\end{lenum}
\end{pro}

\begin{pro}
Can a graph with $6$ vertices have vertices with the following degrees:
$3, 4, 1, 5, 4, 2$?  If so, draw it.  If not, prove it.
\end{pro}

\begin{pro}
Prove or disprove that $Q_n$ is bipartite for $n\geq 1$.
\end{pro}

\begin{pro}
For what values of $n$ is $K_n$ bipartite?
\end{pro}

\begin{pro}
Give the adjacency matrix representation of $Q_3$, numbering the vertices in the obvious order.
\end{pro}

\begin{pro}
\begin{lenum}
\item
Give the adjacency matrix representation for $K_4$.
\item
Give the adjacency list representation for $K_4$.
\end{lenum}
\end{pro}

\begin{pro}
Describe what the adjacency matrix looks like for $K_n$ for $n>1$.
\end{pro}

\begin{pro}
Describe what the adjacency matrix looks like for $C_n$ for $n>1$.
\end{pro}

\begin{pro}
Given an adjacency matrix for $C_n$, with $n>1$, how can you modify it to make it the
adjacency matrix for $P_n$?
\end{pro}

\begin{pro}
What property does the adjacency matrix of every undirected graph have that 
is not necessarily true of directed graphs?
\end{pro}

\begin{pro}
Let $G$ be a graph and let $u$ and $v$ be vertices of $G$.
\begin{lenum}
\item
If $G$ is {\em undirected} and there is a path from $u$ to $v$, 
is there necessarily a path from $v$ to $u$?
Explain, giving an example if possible.
\item
If $G$ is {\em directed} and there is a path from $u$ to $v$, 
is there necessarily a path from $v$ to $u$?
Explain, giving an example if possible.
\end{lenum}
\end{pro}

\begin{pro}
For what values of $n$ is $Q_n$ Eulerian?  Prove your claim.
\end{pro}

\begin{pro}
Is $C_n$ Eulerian for all $n\geq 3$?  Prove it or give a counter example.
\end{pro}

\begin{pro}
Prove that $K_n$ is Hamiltonian for all $n\geq 3$.
\end{pro}

\begin{pro}
Prove that $K_{n,n}$ is Hamiltonian for all $n\geq 3$.
\end{pro}

\begin{pro}
For what values of $m$ and $n$ is $K_{m,n}$ Eulerian?
\end{pro}

\begin{pro}
A graph is Eulerian if and only if its adjacency matrix has what property?
\end{pro}

\begin{pro}
What properties does an adjacency matrix for graph $G$ need in order to 
use Theorem~\ref{thm:dirac} to prove it is Hamiltonian?
\end{pro}

\begin{pro}
Let $G$ be a bipartite graph with $v$ vertices and $e$ edges.
Prove that if $e> 2v-4$, then $G$ is not planar.
\end{pro}

\begin{pro}
For each of the following, either give a planar embedding or prove the graph 
is not planar.
\begin{lenum}
\item $Q_3$
%\item $Q_4$ (Hint: $Q_4$ does not contain $C_3$ as a subgraph.)
\item $Q_5$ 
\item $K_{2,3}$
\item $K_6$
\end{lenum}
\end{pro}

% Now an RQ
%\begin{pro}
%If a graph has very few edges, which is better:  an adjacency matrix or an adjacency list?
%Explain.
%\end{pro}

\begin{pro}\label{pro:impcomp}
Let $G$ be a graph with $n$ vertices and $m$ edges  
and let $u$ and $v$ be arbitrary vertices of $G$.
Describe an algorithm that accomplishes each of the following
assuming $G$ is represented using an {\em adjacency matrix}. 
Then give a tight bound on the worst-case complexity of the algorithm.
Your bounds might be based on $n$, $m$, $deg(u)$, and/or $deg(v)$.
\begin{lenum}
\item
Determine the degree of $u$.
\item
Determine whether or not edge $(u,v)$ is in the graph.
\item
Iterate over the neighbors of $u$ (and doing something for each neighbor,
but don't worry about what and assume it takes constant time for each neighbor).
\item
Add an edge between $u$ and $v$.
\end{lenum}
\end{pro}

\begin{pro}
Repeat Problem~\ref{pro:impcomp}, but this time assume that
$G$ is represented using {\em adjacency lists}. 
\begin{lenum}
\item
Determine the degree of $u$.
\item
Determine whether or not edge $(u,v)$ is in the graph.
\item
Iterate over the neighbors of $u$ (and doing something for each neighbor,
but don't worry about what).
\item
Add an edge between $u$ and $v$.
\end{lenum}
\end{pro}

\begin{pro}
\begin{lenum}
\item
List several advantages that the adjacency matrix representation has over the adjacency list representation.
\item
List several advantages that the adjacency list representation has over the adjacency matrix representation.
\end{lenum}
\end{pro}

